% !TeX TXS-program:bibliography = txs:///biber
% chktex-file 1
\documentclass[fontsize=14pt, russian]{scrartcl}
\input{header.tex}
\addbibresource{biblio.bib}

\begin{document}
\sloppy

\def\figurename{Рисунок}

\newcommand{\sqlkw}[1]{\texttt{#1}}
\newcommand{\SELECT}{\sqlkw{SELECT}}
\newcommand{\FROM}{\sqlkw{FROM}}
\newcommand{\WHERE}{\sqlkw{WHERE}}
\newcommand{\JOIN}{\sqlkw{JOIN}}
\newcommand{\GROUPBY}{\sqlkw{GROUP BY}}
\newcommand{\ORDERBY}{\sqlkw{ORDER BY}}

\newcommand{\SeqScan}{\sqlkw{SeqScan}}
\newcommand{\HashJoin}{\sqlkw{HashJoin}}
\newcommand{\NestedLoopJoin}{\sqlkw{NestedLoopJoin}}
\newcommand{\Sort}{\sqlkw{Sort}}
\newcommand{\FilterOp}{\sqlkw{Filter}}
\newcommand{\ProjectionOp}{\sqlkw{Projection}}

\newcommand{\screenplaceholder}[3]{%
  \begin{figure}[H]
    \centering
    \fbox{%
      \begin{minipage}[c][2.5cm][c]{0.91\textwidth}
        \centering
        \textbf{МЕСТО ДЛЯ СКРИНШОТА}\\[0.4cm]
        #3
      \end{minipage}%
    }
    \caption{#1}
    \label{#2}
  \end{figure}
}

\setcounter{page}{2}

\newpage
\renewcommand\contentsname{\hfill{\normalfont{СОДЕРЖАНИЕ}}\hfill}
\tableofcontents
\newpage

\anonsection{ВВЕДЕНИЕ}

Одним из определяющих факторов конкурентоспособности между различными
реализациями СУБД является качество оптимизатора запросов. Оптимизатор
представляет собой компонент, ответственный за преобразование декларативного
SQL-запроса в эффективный физический план исполнения. Именно от качества работы
оптимизатора в значительной степени зависит производительность всей системы:
выбор неоптимального плана исполнения может привести к увеличению времени
выполнения запроса в десятки раз относительно оптимального случая.

Также оптимизатор запросов считается наиболее сложной составной частью любой
СУБД. Задача нахождения оптимального плана является NP-полной в общем случае,
что обуславливает необходимость применения эвристических методов и алгоритмов
ограниченного перебора. Одним из первых практических подходов к стоимостной
оптимизации стал алгоритм System~R~\cite{Selinger1979}. В дальнейшем были
разработаны расширяемые архитектуры Volcano~\cite{GraefeMcKenna1993} и
Cascades~\cite{Graefe1995}.

Преимущества архитектуры Cascades заключаются в расширяемости, модульности и
возможности управлять поиском с помощью набора правил.

Целью практики является разработка и исследование программного прототипа
SQL-компилятора с оптимизатором запросов на основе архитектуры Cascades.

Для достижения поставленной цели необходимо решить следующие задачи:
\begin{itemize}
  \item изучить существующие архитектуры оптимизаторов запросов;
  \item реализовать разбор подмножества SQL и внутреннее представление
        операторов реляционной алгебры;
  \item реализовать структуру данных Memo и алгоритм поиска оптимального плана
        на основе архитектуры Cascades;
  \item реализовать набор правил трансформации и реализации для основных
        операторов реляционной алгебры;
  \item реализовать метод ветвей и границ для эффективного отсечения
        неоптимальных планов;
  \item реализовать исполнение физических планов над таблицами в формате CSV;
  \item добавить интерпретацию и JIT-компиляцию выражений;
  \item подготовить тесты и вспомогательные средства для исследования планов.
\end{itemize}

В рамках практики разработана модельная СУБД, включающая основные этапы
обработки SQL-запроса. Программа разбирает текст запроса, строит дерево
операторов реляционной алгебры, выполняет стоимостную оптимизацию, формирует
физический план и исполняет его над таблицами в формате CSV. Дополнительно
реализована JIT-компиляция скалярных выражений с использованием LLVM.

\section{ПРОЕКТИРОВАНИЕ}

\subsection{Назначение и область применения}

Разработанная программа является исследовательским прототипом обработчика
SQL-запросов. Она не заменяет промышленную СУБД, поскольку не содержит
транзакционного менеджера, журналирования, многоверсионности и сетевого
протокола. Программа предназначена для изучения этапов компиляции запросов,
экспериментальной проверки правил оптимизации и измерения стоимости физических
операторов.

Входными данными служат SQL-запрос и директория с CSV-файлами. Каждый файл
соответствует одной таблице. Первая строка файла задает схему отношения, а
остальные строки содержат кортежи. Результат запроса выводится в стандартный
поток вывода в текстовом виде.

Поддерживаются запросы \SELECT{} с предложениями \FROM{} и \WHERE{},
проекции столбцов и выражений, соединения и декартовы произведения таблиц.
Скалярные выражения включают сравнения, \sqlkw{BETWEEN}, \sqlkw{IN}, строковые
и целочисленные значения, а также \sqlkw{NULL}. Реализованы агрегатные функции
\sqlkw{SUM}, \sqlkw{COUNT}, группировка \GROUPBY{} и сортировка \ORDERBY{}.

Некоторые возможности полного SQL намеренно не реализованы. К ним относятся
\sqlkw{DISTINCT}, \sqlkw{HAVING}, оконные функции, подзапросы и операции над
множествами. При встрече неподдерживаемой конструкции парсер возвращает
диагностическое сообщение.

\subsection{Общая архитектура программы}

Сначала программа читает SQL-запрос, выполняет лексический и синтаксический
анализ, затем строит логическое представление. Оптимизатор заполняет Memo,
перечисляет эквивалентные выражения и выбирает физический план с минимальной
стоимостью. Исполнитель обрабатывает план над CSV-файлами и выводит схему и
кортежи результирующего отношения.

\section{РЕАЛИЗАЦИЯ}

\subsection{Выбор инструментов разработки}

Программа написана на языке C++ с использованием стандарта C++23. Этот язык
позволяет контролировать представление данных в памяти, использовать
стандартные алгебраические типы и корутины, а также интегрироваться с LLVM.
Для сборки применяется CMake. Воспроизводимое окружение разработки описано
в файле \texttt{flake.nix}.

Для синтаксического анализа выбран генератор парсеров ANTLR4. Асинхронное
взаимодействие физических операторов построено на Boost.Asio. LLVM
используется для генерации машинного кода скалярных выражений. Модульные тесты
реализованы с помощью GoogleTest, а измерения производительности физических
операторов --- с помощью Google Benchmark.

\subsection{Модуль синтаксического анализа}

Для синтаксического анализа используется грамматика PostgreSQL и генератор
парсеров ANTLR4. Сгенерированные классы лексического и синтаксического
анализаторов находятся в каталоге
\texttt{src/stewkk/sql/logic/parser/codegen}. Функция \texttt{GetAST}
инициализирует анализаторы, строит дерево разбора и запускает обход дерева.

Обход реализован классом \texttt{Visitor}. Для интересующих конструкций SQL
переопределены методы посещения узлов грамматики. Например, предложение
\WHERE{} преобразуется в логический оператор фильтрации, список выражений
после \SELECT{} преобразуется в проекцию, а список таблиц после \FROM{} ---
в дерево соединений.

Результатом работы парсера является структура \texttt{ParsedQuery}. Она
содержит логический оператор верхнего уровня и необязательное требование к
порядку строк для \ORDERBY{}. Логические операторы и скалярные выражения
представлены алгебраическими типами на основе \texttt{std::variant}. В них
входят операторы таблицы, фильтрации, проекции, агрегации, соединения, ссылки
на атрибуты, константы и арифметические выражения.

\subsection{Модуль оптимизации запросов}

Оптимизатор получает логическое дерево, набор правил, оценки кардинальности,
каталог схем и требуемые свойства результата. Центральной структурой данных
является Memo. Она хранит группы логически эквивалентных выражений.

Каждая группа содержит логические и физические выражения. Дочерние узлы
ссылаются на группы, поэтому одно выражение компактно представляет множество
возможных поддеревьев. Повторяющиеся выражения не добавляются в Memo повторно.

Правила оптимизации разделены на две категории. Трансформационные правила
создают новые логические выражения, эквивалентные исходным. Реализационные
правила добавляют физические операторы, которые могут быть непосредственно
исполнены.

В текущей версии реализованы следующие трансформации:
\begin{itemize}
  \item коммутативность соединения;
  \item ассоциативность соединения;
  \item разбиение фильтра с конъюнкцией на последовательность фильтров;
  \item объединение последовательных фильтров;
  \item проталкивание фильтра через проекцию;
  \item проталкивание фильтра во входы соединения;
  \item преобразование списка \sqlkw{IN} в цепочку сравнений.
\end{itemize}

Реализационные правила создают последовательное сканирование, фильтрацию,
проекцию, агрегацию, соединение вложенными циклами, декартово произведение
вложенными циклами и хеш-соединение. Хеш-соединение применяется только для
простого внутреннего соединения по равенству двух атрибутов.

Поиск выполняется сверху вниз с использованием стека отложенных задач. Для
каждой пары из группы и требуемого набора физических свойств сохраняется
лучший найденный вариант. Такой вариант называется победителем группы.
Стоимость полного плана равна сумме локальных стоимостей его операторов.

Для уменьшения пространства поиска используется метод ветвей и границ.
Оптимизатор вычисляет допустимую нижнюю оценку стоимости группы. Если даже
нижняя оценка не меньше стоимости уже известного решения, дальнейшее
исследование ветви прекращается.

Физические свойства описывают дополнительные требования к результату. В
текущей версии реализовано свойство сортировки. Если запрос содержит
\ORDERBY{}, оптимизатор должен построить план, выдающий строки в требуемом
порядке. При необходимости в план добавляется обеспечивающий оператор
\Sort{}.

\subsection{Физический план и его сериализация}

После завершения поиска оптимизатор преобразует выбранные физические выражения
в дерево \texttt{PhysicalPlanNode}. Физический план отделен от внутренней
структуры Memo и может быть исполнен независимо от процесса оптимизации.

К основным физическим операторам относятся последовательное чтение CSV-файла
\SeqScan{}, фильтрация \FilterOp{}, проекция \ProjectionOp{}, соединения
\NestedLoopJoin{} и \HashJoin{}, сортировка \Sort{} и хеш-агрегация
\sqlkw{HashAggregate}.

Планы сериализуются в текстовый формат S-expression. Его удобно читать,
сравнивать в тестах и передавать вспомогательным утилитам. Дополнительно
реализован экспорт плана в формат Graphviz DOT. После запуска программы файл
сохраняется в каталоге \texttt{.plans}.

\subsection{Модуль хранения и исполнения данных}

Каждой таблице соответствует CSV-файл с совпадающим именем. Например, таблица
\texttt{users} хранится в файле \texttt{users.csv}. Первая строка содержит
описание столбцов и их типов:

\begin{Verbatim}[fontsize=\small]
id:int,name:string,age:int
1,Alice,25
2,Bob,31
\end{Verbatim}

Исполнитель физических планов реализован классом \texttt{Executor}. Он
параметризуется способом вычисления выражений и функциями доступа к таблицам.
Благодаря этому один исполнитель используется с интерпретатором, JIT-компилятором
и различными реализациями сканирования.

Операторы взаимодействуют через каналы Boost.Asio. Отдельные каналы передают
описание атрибутов и порции размером до 2048 кортежей. Это уменьшает число
операций синхронизации по сравнению с передачей отдельных строк.

Проекция и фильтрация являются потоковыми операторами. Они получают порцию
входных кортежей, обрабатывают ее и передают результат следующему оператору.
Для блокирующих операций может потребоваться накопление данных. Например,
сортировка сохраняет все входные кортежи в памяти.

Соединение вложенными циклами материализует один вход во временный бинарный
файл, после чего многократно читает его при обработке второго входа.
Хеш-соединение строит хеш-таблицу по ключу левого входа и выполняет поиск
совпадений для кортежей правого входа.

Для строковых значений используется интернирование. Строка помещается в общий
пул, а внутри кортежа хранится целочисленный идентификатор. Это сохраняет
компактное представление значения и упрощает передачу кортежей между
операторами.

\subsection{JIT-компиляция выражений}

Для часто вычисляемых выражений реализована JIT-компиляция с использованием
LLVM ORC JIT. Класс \texttt{JITCompiler} преобразует дерево выражения в LLVM IR,
запускает оптимизирующие проходы и возвращает указатель на машинную функцию.

Скомпилированная функция получает указатель на результирующее значение,
указатель на входной кортеж и описание атрибутов. Структура значения содержит
флаг \texttt{is\_null} и объединение для непосредственного значения. Такое
представление учитывает трехзначную логику SQL.

В проекте доступны интерпретируемая, JIT-компилируемая и кэшируемая
JIT-компилируемая стратегии вычисления выражений. Последняя сохраняет
полученные машинные функции и переиспользует их.

Текущая реализация JIT предназначена для целочисленных и логических выражений.
Строковые операции, \sqlkw{IN}, агрегатные выражения и возведение в степень
через JIT не поддерживаются. Для них следует использовать интерпретируемый
режим.

\subsection{Руководство пользователя}

\subsubsection{Требования к окружению}

Для сборки программы требуются CMake, компилятор C++ с поддержкой C++23, LLVM,
ANTLR4 и библиотеки Boost. В репозитории присутствует файл \texttt{flake.nix},
поэтому в системе с пакетным менеджером Nix рекомендуется открыть окружение:

\begin{Verbatim}[fontsize=\small]
nix develop
\end{Verbatim}

Все команды далее выполняются из корневого каталога проекта.

\subsubsection{Сборка программы}

Для первичной конфигурации и сборки следует выполнить:

\begin{Verbatim}[fontsize=\small]
cmake -S . -B build -DCMAKE_BUILD_TYPE=Release
cmake --build build -- -j 6
\end{Verbatim}

После успешной сборки исполняемый файл находится по пути
\texttt{build/bin/sql}. Для повторной сборки также можно использовать цель
\texttt{make build}.

\subsubsection{Подготовка входных данных}

Пользователь должен создать отдельную директорию и поместить в нее CSV-файлы.
Имя каждого файла без расширения считается именем таблицы. В заголовке каждого
файла для столбца указывается имя и тип через двоеточие. Поддерживаются типы
\texttt{int} и \texttt{string}. Пустое SQL-значение записывается как
\texttt{NULL}.

Для первого запуска можно использовать готовые тестовые данные:

\begin{Verbatim}[fontsize=\small]
test/static/executor/test_data
\end{Verbatim}

\subsubsection{Запуск SQL-запроса}

SQL-запрос передается программе через стандартный поток ввода. Простейший
пример запуска:

\begin{Verbatim}[fontsize=\small]
echo 'SELECT users.id FROM users;' | \
  ./build/bin/sql --data-dir test/static/executor/test_data
\end{Verbatim}

Для запуска запроса из файла следует использовать перенаправление ввода:

\begin{Verbatim}[fontsize=\small]
./build/bin/sql --data-dir path/to/csv_dir < query.sql
\end{Verbatim}

Программа выводит заголовок отношения, содержащий имена и типы столбцов, а
затем полученные строки. Если запрос не содержит \ORDERBY{}, строки вывода
сортируются лексикографически для воспроизводимости результата.

\screenplaceholder{Выполнение SQL-запроса над тестовыми данными}{fig:practice-run}{%
  Вставить скриншот запуска запроса через стандартный поток ввода\\
  и напечатанного результирующего отношения.}

\subsubsection{Вывод логического и физического планов}

Для диагностики предусмотрены параметры \texttt{--print-ast} и
\texttt{--print-plan}. Первый параметр выводит логическое дерево после
синтаксического анализа, второй --- выбранный физический план:

\begin{Verbatim}[fontsize=\small]
echo 'SELECT users.id FROM users WHERE users.age > 18 \
ORDER BY users.id;' | \
  ./build/bin/sql --data-dir test/static/executor/test_data \
  --print-ast --print-plan
\end{Verbatim}

\screenplaceholder{Вывод логического дерева и физического плана}{fig:practice-plan}{%
  Вставить скриншот запуска с параметрами\\
  \texttt{--print-ast --print-plan}.}

\subsubsection{Запуск с JIT-компиляцией}

Для выбора JIT-компиляции выражений используется параметр \texttt{--jit}:

\begin{Verbatim}[fontsize=\small]
echo 'SELECT users.id FROM users WHERE users.age > 18;' | \
  ./build/bin/sql --data-dir test/static/executor/test_data --jit
\end{Verbatim}

Параметр рекомендуется применять к запросам с целочисленными арифметическими и
логическими выражениями. Запросы со строковыми выражениями и \sqlkw{IN} следует
запускать без JIT.

\screenplaceholder{Выполнение запроса с JIT-компиляцией выражений}{fig:practice-jit}{%
  Вставить скриншот запуска программы с параметром \texttt{--jit}.}

\subsubsection{Доступные параметры командной строки}

\begin{longtable}{|p{4.2cm}|p{9.3cm}|}
  \caption{Параметры командной строки программы.}\label{tbl:practice-cli}\\
  \hline
  Параметр & Назначение \\
  \hline
  \endfirsthead
  \multicolumn{2}{l}{\small Продолжение таблицы~\ref{tbl:practice-cli}}\\
  \hline
  Параметр & Назначение \\
  \hline
  \endhead
  \texttt{--data-dir DIR} &
  Директория с CSV-файлами таблиц. Обязательна при обычном исполнении запроса. \\
  \hline
  \texttt{--print-ast} &
  Вывод логического дерева операторов после синтаксического анализа. \\
  \hline
  \texttt{--print-plan} &
  Вывод выбранного физического плана в формате S-expression. \\
  \hline
  \texttt{--jit} &
  Использование кэшируемой JIT-компиляции скалярных выражений. \\
  \hline
  \texttt{--check-reachable FILE} &
  Проверка достижимости физического плана, сериализованного в файле. \\
  \hline
\end{longtable}

\section{ТЕСТИРОВАНИЕ}

Модульные тесты написаны с использованием GoogleTest. Они проверяют парсер,
вычисление выражений, исполнение операторов, сериализацию планов, свойства
сортировки, Memo и оптимизатор. Отдельные тесты проверяют интерпретируемое и
JIT-исполнение.

Для проверки корректности C++-части проекта следует выполнить:

\begin{Verbatim}[fontsize=\small]
ctest --test-dir build --output-on-failure
\end{Verbatim}

В результате запуска успешно выполняются 126 тестов. Для проверки конвертера
запросов Star Schema Benchmark используется команда:

\begin{Verbatim}[fontsize=\small]
pytest -q benchmarks/datasets/ssb
\end{Verbatim}

Набор содержит 17 тестов, которые также выполняются успешно. В каталоге
\texttt{benchmarks} находятся тесты производительности на основе Google
Benchmark. Они предназначены для измерения характеристик отдельных физических
операторов.

Вспомогательные Python-скрипты из каталога \texttt{research} генерируют
запросы, сравнивают результаты с Microsoft SQL Server и проверяют достижимость
внешних физических планов. Для последней задачи программа поддерживает режим
\texttt{--check-reachable}, запускающий полный перебор пространства планов без
отсечения.

\anonsection{ЗАКЛЮЧЕНИЕ}

В ходе практики разработан программный прототип SQL-компилятора и модельной
системы исполнения запросов. Реализован сквозной процесс от разбора SQL-текста
до получения результирующего отношения. Программа работает с CSV-файлами,
поддерживает основные операторы реляционной алгебры и формирует физический план
на основе стоимостной модели.

Ключевой частью проекта является расширяемый оптимизатор, использующий
структуру Memo и правила преобразования. Реализован поиск планов с учетом
физических свойств, нижних оценок стоимости и метода ветвей и границ.
Поддержка обеспечивающего оператора сортировки позволяет корректно обрабатывать
требования \ORDERBY{}.

Исполнитель построен на корутинах C++ и каналах Boost.Asio. Потоковая передача
порций кортежей сочетается с материализацией для блокирующих операторов.
Для ускорения вычисления выражений добавлена JIT-компиляция с использованием
LLVM и кэширование скомпилированных функций.

Корректность основных компонентов проверяется модульными тестами. Дополнительно
подготовлены бенчмарки физических операторов и инструменты дифференциального
сравнения с Microsoft SQL Server. Эти средства позволяют использовать проект
как основу для дальнейшего исследования правил оптимизации и уточнения
стоимостной модели.

\renewcommand\refname{СПИСОК ИСПОЛЬЗОВАННЫХ ИСТОЧНИКОВ}
{\catcode`"\active\def"{\relax}
    \addcontentsline{toc}{section}{\protect\numberline{}\refname}%
    \printbibliography
}

\end{document}
